\documentclass[11pt,a4paper]{article}
\usepackage[utf8]{inputenc}
\usepackage[vietnamese]{babel}
\usepackage[T5]{fontenc}
\usepackage{amsmath,amssymb,amsfonts}
\usepackage{geometry}
\usepackage{fancyhdr}
\usepackage{multicol}
\usepackage{xcolor}
\usepackage{tikz}
\usetikzlibrary{calc}

% Thiết lập khổ giấy A4 và căn lề tối ưu
\geometry{
	a4paper,
	left=15mm,
	right=15mm,
	top=15mm,
	bottom=18mm
}

% Cấu hình header & footer
\pagestyle{fancy}
\fancyhf{}
\renewcommand{\headrulewidth}{0pt}
\renewcommand{\footrulewidth}{0.4pt}
\rfoot{\footnotesize\color{subgray}Trang \thepage}
\lfoot{\footnotesize\color{subgray}Chuyên đề Toán 10 -- Mệnh đề \& Tập hợp (GDPT 2018)}

% Định nghĩa bảng màu Infographic học thuật
\definecolor{primaryred}{HTML}{C62828}   % Đỏ son cho tiêu đề
\definecolor{inkblue}{HTML}{182A78}      % Xanh mực cho đường nét & khung
\definecolor{subgray}{HTML}{424242}      % Xám cho thông tin phụ

% ------------------------------------------------------------------------------
% ĐỊNH NGHĨA CÁC MÔI TRƯỜNG & LỆNH TƯƠNG THÍCH CHUẨN GÓI ex_test
% ------------------------------------------------------------------------------
\newcounter{ex}
\newenvironment{ex}[1][]{%
	\refstepcounter{ex}%
	\par\smallskip\noindent\textbf{\color{inkblue}Câu \theex.} #1\ignorespaces
}{\par}

% \True là cờ đánh dấu phương án đúng (chuẩn ex_test)
\newcommand{\True}{}
\renewcommand{\True}{\color{red}\textbf{\underline{X}} }

% Lệnh \choice tự động điều chỉnh 1 dòng, 2 dòng hoặc 4 dòng tuỳ độ dài phương án
\newlength{\widthcha}
\newlength{\widthchb}
\newlength{\widthchc}
\newlength{\widthchd}

\newcommand{\choiceFour}[4]{%
	\par\smallskip
	\noindent\textbf{A.} #1\par
	\noindent\textbf{B.} #2\par
	\noindent\textbf{C.} #3\par
	\noindent\textbf{D.} #4\par\smallskip
}

\newcommand{\choiceTwo}[4]{%
	\ifdim\widthcha<0.46\linewidth
	\ifdim\widthchb<0.46\linewidth
	\ifdim\widthchc<0.46\linewidth
	\ifdim\widthchd<0.46\linewidth
	\par\smallskip\noindent
	\begin{minipage}[t]{0.48\linewidth}\textbf{A.} #1\end{minipage}\hfill
	\begin{minipage}[t]{0.48\linewidth}\textbf{B.} #2\end{minipage}\par\smallskip\noindent
	\begin{minipage}[t]{0.48\linewidth}\textbf{C.} #3\end{minipage}\hfill
	\begin{minipage}[t]{0.48\linewidth}\textbf{D.} #4\end{minipage}\par\smallskip
	\else\choiceFour{#1}{#2}{#3}{#4}\fi
	\else\choiceFour{#1}{#2}{#3}{#4}\fi
	\else\choiceFour{#1}{#2}{#3}{#4}\fi
	\else\choiceFour{#1}{#2}{#3}{#4}\fi
}

\newcommand{\choice}[4]{%
	\settowidth{\widthcha}{\textbf{A.} #1}%
	\settowidth{\widthchb}{\textbf{B.} #2}%
	\settowidth{\widthchc}{\textbf{C.} #3}%
	\settowidth{\widthchd}{\textbf{D.} #4}%
	\ifdim\widthcha<0.22\linewidth
	\ifdim\widthchb<0.22\linewidth
	\ifdim\widthchc<0.22\linewidth
	\ifdim\widthchd<0.22\linewidth
	\par\smallskip\noindent
	\begin{minipage}[t]{0.24\linewidth}\textbf{A.} #1\end{minipage}\hfill
	\begin{minipage}[t]{0.24\linewidth}\textbf{B.} #2\end{minipage}\hfill
	\begin{minipage}[t]{0.24\linewidth}\textbf{C.} #3\end{minipage}\hfill
	\begin{minipage}[t]{0.24\linewidth}\textbf{D.} #4\end{minipage}\par\smallskip
	\else\choiceTwo{#1}{#2}{#3}{#4}\fi
	\else\choiceTwo{#1}{#2}{#3}{#4}\fi
	\else\choiceTwo{#1}{#2}{#3}{#4}\fi
	\else\choiceTwo{#1}{#2}{#3}{#4}\fi
}

% Môi trường \choiceTF cho câu hỏi trắc nghiệm Đúng/Sai
\newcommand{\choiceTF}[4]{%
	\par\smallskip
	\noindent\begin{minipage}[t]{0.82\linewidth}\textbf{a)} #1\end{minipage}\hfill\fbox{\small\strut\ \textbf{Đ}\ \ \textbar\ \ \textbf{S}\ }\par\smallskip
	\noindent\begin{minipage}[t]{0.82\linewidth}\textbf{b)} #2\end{minipage}\hfill\fbox{\small\strut\ \textbf{Đ}\ \ \textbar\ \ \textbf{S}\ }\par\smallskip
	\noindent\begin{minipage}[t]{0.82\linewidth}\textbf{c)} #3\end{minipage}\hfill\fbox{\small\strut\ \textbf{Đ}\ \ \textbar\ \ \textbf{S}\ }\par\smallskip
	\noindent\begin{minipage}[t]{0.82\linewidth}\textbf{d)} #4\end{minipage}\hfill\fbox{\small\strut\ \textbf{Đ}\ \ \textbar\ \ \textbf{S}\ }\par\smallskip
}

% Lệnh điền đáp số cho phần trả lời ngắn
\newcommand{\shortans}[2][]{%
	\par\smallskip\noindent\textbf{Đáp số:} \dotfill\ \framebox[2.4cm]{\rule{0pt}{1.1em}#2}\par
}

% Định dạng tiêu đề các phân phần
\newcommand{\sectionheader}[2]{%
	\vspace{0.35cm}
	\noindent
	\begin{tikzpicture}
		\node[fill=primaryred!10, draw=primaryred, rounded corners=3pt, inner sep=4pt, text width=\linewidth-8pt] {
			\textbf{\color{primaryred}#1}\\[1pt]
			\footnotesize\textit{\color{subgray}#2}
		};
	\end{tikzpicture}
	\par\vspace{0.15cm}
}

\begin{document}
	\renewcommand{\shortans}[2][]{\par\smallskip\noindent\textbf{Đáp số:} \dotfill\ \framebox[2.4cm]{\rule{0pt}{1.1em}\color{red}\textbf{#2}}\par}
	
	% ------------------------------------------------------------------------------
	% TIÊU ĐỀ BÀI HỌC VÀ THÔNG TIN HỌC SINH
	% ------------------------------------------------------------------------------
	\begin{center}
		{\large \textbf{\color{primaryred}PHIẾU HỌC TẬP INFOGRAPHIC: TẬP HỢP VÀ CÁC PHÉP TOÁN TẬP HỢP}}\\[1.5mm]
		{\small \textbf{\color{inkblue}Chuyên đề Toán 10 -- Chương I: Mệnh đề và Tập hợp (GDPT 2018)}}\\[2mm]
		\noindent
		\begin{tikzpicture}
			\draw[color=inkblue!40, line width=0.8pt] (0,0) -- (\linewidth,0);
		\end{tikzpicture}\\[1mm]
		{\footnotesize \textbf{Họ và tên:} \dots\dots\dots\dots\dots\dots\dots\dots\dots\dots\dots\dots\dots\dots \quad 
			\textbf{Lớp:} \dots\dots\dots\dots\dots \quad 
			\textbf{Ngày:} \dots\dots/\dots\dots/20\dots\dots}
	\end{center}
	
	% ==============================================================================
	% PHẦN I: TRẮC NGHIỆM NHIỀU LỰA CHỌN (10 CÂU - CHIA 2 CỘT)
	% ==============================================================================
	\sectionheader{PHẦN I. TRẮC NGHIỆM NHIỀU LỰA CHỌN (10 CÂU)}{Khoanh tròn vào chữ cái A, B, C hoặc D trước phương án trả lời đúng nhất.}
	
	\setlength{\columnsep}{18pt}
	\begin{multicols}{2}
		
		\begin{ex}
			Cho tập hợp $A = \{x \in \mathbb{R} \mid (x^2 - 4)(x - 1) = 0\}$. Khẳng định nào sau đây là đúng?
			\choice
			{\True $A = \{-2, 1, 2\}$}
			{$A = \{1, 2\}$}
			{$A = \{-2, 2\}$}
			{$A = \{-1, 1, 2\}$}
		\end{ex}
		
		\begin{ex}
			Ký hiệu nào sau đây thể hiện đúng phát biểu: ``Số $2$ là một số tự nhiên''?
			\choice
			{\True $2 \in \mathbb{N}$}
			{$2 \subset \mathbb{N}$}
			{$\{2\} \in \mathbb{N}$}
			{$2 = \mathbb{N}$}
		\end{ex}
		
		\begin{ex}
			Cho hai tập hợp $A = \{1, 2, 3, 4, 5\}$ và $B = \{2, 4, 6, 8\}$. Tìm tập hợp $A \cap B$.
			\choice
			{\True $\{2, 4\}$}
			{$\{1, 2, 3, 4, 5, 6, 8\}$}
			{$\{1, 3, 5\}$}
			{$\{6, 8\}$}
		\end{ex}
		
		\begin{ex}
			Cho hai tập hợp $A = \{-1, 0, 1, 2\}$ và $B = \{0, 2, 4\}$. Tìm tập hợp hiệu $A \setminus B$.
			\choice
			{\True $\{-1, 1\}$}
			{$\{0, 2\}$}
			{$\{4\}$}
			{$\{-1, 0, 1, 2, 4\}$}
		\end{ex}
		
		\begin{ex}
			Cho tập hợp $A = (-2; 3]$ và $B = (1; 5)$. Tìm tập hợp $A \cup B$.
			\choice
			{\True $(-2; 5)$}
			{$(1; 3]$}
			{$[-2; 5]$}
			{$(-2; 1)$}
		\end{ex}
		
		\begin{ex}
			Cho tập hợp $A = [-1; 4]$. Tìm phần bù $C_{\mathbb{R}} A$.
			\choice
			{\True $(-\infty; -1) \cup (4; +\infty)$}
			{$(-\infty; -1] \cup [4; +\infty)$}
			{$(-\infty; 4)$}
			{$(4; +\infty)$}
		\end{ex}
		
		\begin{ex}
			Tập hợp $A = \{a, b, c\}$ có tất cả bao nhiêu tập hợp con?
			\choice
			{\True $8$}
			{$6$}
			{$7$}
			{$9$}
		\end{ex}
		
		\begin{ex}
			Cho hai khoảng $A = (-\infty; m)$ và $B = (3; +\infty)$. Tìm điều kiện của tham số $m$ để $A \cap B = \emptyset$.
			\choice
			{\True $m \le 3$}
			{$m < 3$}
			{$m > 3$}
			{$m \ge 3$}
		\end{ex}
		
		\begin{ex}
			Cho hai tập hợp $A = [m-1; m+3]$ và $B = (-2; 5)$. Tìm tất cả các giá trị của $m$ để $A \subset B$.
			\choice
			{\True $-1 < m < 2$}
			{$-1 \le m \le 2$}
			{$m > -1$}
			{$m < 2$}
		\end{ex}
		
		\begin{ex}
			Lớp 10A có 42 học sinh, trong đó có 25 học sinh thích môn Toán, 18 học sinh thích môn Tiếng Anh và 8 học sinh không thích môn nào trong hai môn này. Hỏi có bao nhiêu học sinh thích cả hai môn Toán và Tiếng Anh?
			\choice
			{\True $9$}
			{$7$}
			{$11$}
			{$5$}
		\end{ex}
		
	\end{multicols}
	
	% ==============================================================================
	% PHẦN II: TRẮC NGHIỆM ĐÚNG - SAI (4 CÂU)
	% ==============================================================================
	\sectionheader{PHẦN II. TRẮC NGHIỆM ĐÚNG -- SAI (4 CÂU)}{Học sinh ghi chữ Đ (Đúng) hoặc S (Sai) vào ô trống tương ứng trước mỗi nhận định.}
	
	\begin{ex}
		Cho hai tập hợp $A = \{x \in \mathbb{R} \mid x^2 - 5x + 6 = 0\}$ và $B = \{x \in \mathbb{Z} \mid -1 \le x \le 3\}$:
		\choiceTF
		{\True Viết tập hợp $A$ dưới dạng liệt kê phần tử ta được $A = \{2, 3\}$.}
		{Số phần tử của tập hợp $B$ là $4$ phần tử.}
		{\True Giao của hai tập hợp là $A \cap B = \{2, 3\}$.}
		{Tập hợp $A$ không phải là tập con của tập hợp $B$ ($A \not\subset B$).}
	\end{ex}
	
	\begin{ex}
		Cho hai tập hợp $A = [-3; 2)$ và $B = (0; 5]$ trên tập số thực $\mathbb{R}$:
		\choiceTF
		{\True Hợp của hai tập hợp là $A \cup B = [-3; 5]$.}
		{\True Giao của hai tập hợp là $A \cap B = (0; 2)$.}
		{\True Hiệu của hai tập hợp là $A \setminus B = [-3; 0]$.}
		{\True Phần bù của $A$ trong $\mathbb{R}$ là $C_{\mathbb{R}} A = (-\infty; -3) \cup [2; +\infty)$.}
	\end{ex}
	
	\begin{ex}
		Xét quan hệ bao hàm giữa các tập hợp số $\mathbb{N}, \mathbb{Z}, \mathbb{Q}, \mathbb{R}$ và tập số vô tỉ $\mathbb{I}$:
		\choiceTF
		{\True Tập hợp các số tự nhiên $\mathbb{N}$ là tập con của tập hợp các số nguyên $\mathbb{Z}$ ($\mathbb{N} \subset \mathbb{Z}$).}
		{\True Phép giao $\mathbb{Z} \cap \mathbb{Q} = \mathbb{Z}$.}
		{\True Phép hiệu $\mathbb{R} \setminus \mathbb{Q}$ là tập hợp các số vô tỉ $\mathbb{I}$.}
		{Phần bù $C_{\mathbb{R}} \mathbb{N} = (0; +\infty)$.}
	\end{ex}
	
	\begin{ex}
		Cho hai tập hợp phụ thuộc tham số $m$: $A = (1; m+1]$ và $B = (2; 6)$:
		\choiceTF
		{\True Điều kiện để tập hợp $A$ khác rỗng ($A \neq \emptyset$) là $m > 0$.}
		{Khi $m = 4$, tập hợp $A = (1; 5]$ và $A \cap B = (2; 5]$.}
		{\True Điều kiện của $m$ để $A \cap B = \emptyset$ là $m \le 1$.}
		{\True Điều kiện của $m$ để $B \subset A$ là $m \ge 5$.}
	\end{ex}
	
	% ==============================================================================
	% PHẦN III: CÂU HỎI TRẢ LỜI NGẮN (5 CÂU)
	% ==============================================================================
	\sectionheader{PHẦN III. CÂU HỎI TRẢ LỜI NGẮN (5 CÂU)}{Học sinh viết đáp số hoặc câu trả lời ngắn gọn vào khoảng trống bên dưới mỗi câu hỏi.}
	
	\begin{ex}
		Cho tập hợp $A = \{x \in \mathbb{R} \mid (x^2 - 9)(x^2 - 4x + 3) = 0\}$. Hỏi tập hợp $A$ có tất cả bao nhiêu tập hợp con?
		\shortans{$8$}
	\end{ex}
	
	\begin{ex}
		Cho hai khoảng $A = [-2; 4]$ và $B = (m; m+3)$. Tìm tất cả các giá trị thực của tham số $m$ để $B \subset A$.
		\shortans{$-2 \le m \le 1$}
	\end{ex}
	
	\begin{ex}
		Cho hai tập hợp $A = (-\infty; m+1]$ và $B = (2; +\infty)$. Tìm tất cả các giá trị thực của tham số $m$ để $A \cap B = \emptyset$.
		\shortans{$m \le 1$}
	\end{ex}
	
	\begin{ex}
		Một công ty phỏng vấn 100 khách hàng về việc sử dụng hai sản phẩm A và B. Kết quả có 65 người dùng sản phẩm A, 45 người dùng sản phẩm B và 15 người không dùng sản phẩm nào. Hỏi có bao nhiêu người dùng CẢ HAI sản phẩm A và B?
		\shortans{$25$}
	\end{ex}
	
	\begin{ex}
		\textbf{(Vận dụng thực tế)} Lớp 10B có 40 học sinh. Trong kỳ thi học sinh giỏi cấp trường, có 22 học sinh thi môn Toán, 18 học sinh thi môn Vật lý và 10 học sinh không thi môn nào. Trường quyết định trao thưởng cho các học sinh thi CHỈ MỘT môn. Hỏi có bao nhiêu học sinh lớp 10B được trao thưởng?
		\shortans{$20$}
	\end{ex}
	
\end{document}
